\documentclass{standalone}
\usepackage{tikz}
\tikzset{>=stealth}

\usepackage{xeCJK}
\setCJKmainfont{Microsoft YaHei}
\special{dvipdfmx:config z 0} %取消PDF压缩，加快速度，最终版本生成的时候最好把这句话注释掉

\begin{document}
    \begin{tikzpicture}
        \draw[rounded corners,fill=green!10,shift ={(0 ,0.3)}] (0,0) rectangle (10,-1) node [midway] {辣椒果胶-壳聚糖复合膜的制备};
            \draw[dashed,fill=yellow!10,shift ={(2 ,-0.85)}] (0,0) rectangle (6,-0.8) node [midway,font=\fontsize{8}{9}\selectfont] {不同分子量壳聚糖涂膜液的制备及筛选};
            \draw[dashed,fill=yellow!10,shift ={(2 ,-1.55)}] (0,0) rectangle (6,-0.8) node [midway,font=\fontsize{8}{9}\selectfont] {辣椒果胶-壳聚糖复合膜的制备};

        \draw[->,shift ={(0 ,-1.5)}] (5,-1) -- (5,-1.5);
        \draw[rounded corners,fill=green!10,shift ={(0 ,-3)}] (0,0) rectangle (10,-1) node [midway] {辣椒果胶-壳聚糖复合膜的结构分析};

            \draw[dashed,fill=yellow!10,shift ={(0.5 ,-4.25)}] (0,0) rectangle (1,-3) node [midway,text width =1em, font=\fontsize{8}{9}\selectfont] {厚度和密度};
            \draw[dashed,fill=yellow!10,shift ={(1.8 ,-4.25)}] (0,0) rectangle (1,-3) node [midway,text width =1em, font=\fontsize{8}{9}\selectfont] {色值和透明度};
            \draw[dashed,fill=yellow!10,shift ={(3.1 ,-4.25)}] (0,0) rectangle (1,-3) node [midway,text width =1em, font=\fontsize{8}{9}\selectfont] {红外光谱检测};
            \draw[dashed,fill=yellow!10,shift ={(4.4 ,-4.25)}] (0,0) rectangle (1,-3) node [midway,text width =1em, font=\fontsize{8}{9}\selectfont] {核磁共振检测};
            \draw[dashed,fill=yellow!10,shift ={(5.7 ,-4.25)}] (0,0) rectangle (1,-3) node [midway,text width =1em, font=\fontsize{8}{9}\selectfont] {X射线衍射检测};
            \draw[dashed,fill=yellow!10,shift ={(7 ,-4.25)}] (0,0) rectangle (1,-3) node [midway,text width =1em, font=\fontsize{8}{9}\selectfont] {差示扫描量热仪检测};
            \draw[dashed,fill=yellow!10,shift ={(8.3 ,-4.25)}] (0,0) rectangle (1,-3) node [midway,text width =1em, font=\fontsize{8}{9}\selectfont] {扫描电镜分析};

        \draw[->,shift ={(0 ,-6.5)}] (5,-1) -- (5,-1.5);
        \draw[rounded corners,fill=green!10,shift ={(0 ,-8)}] (0,0) rectangle (10,-1) node [midway] {辣椒果胶-壳聚糖复合膜的性能测试};
            
            \draw[dashed,fill=yellow!10,shift ={(1.8 ,-9.25)}] (0,0) rectangle (1,-3) node [midway,text width =1em, font=\fontsize{8}{9}\selectfont] {溶解性、溶胀度、含水量};
            \draw[dashed,fill=yellow!10,shift ={(3.1 ,-9.25)}] (0,0) rectangle (1,-3) node [midway,text width =1em, font=\fontsize{8}{9}\selectfont] {水蒸气透过系数};
            \draw[dashed,fill=yellow!10,shift ={(4.4 ,-9.25)}] (0,0) rectangle (1,-3) node [midway,text width =1em, font=\fontsize{8}{9}\selectfont] {机械性能};
            \draw[dashed,fill=yellow!10,shift ={(5.7 ,-9.25)}] (0,0) rectangle (1,-3) node [midway,text width =1em, font=\fontsize{8}{9}\selectfont] {抗氧化性};
            \draw[dashed,fill=yellow!10,shift ={(7 ,-9.25)}] (0,0) rectangle (1,-3) node [midway,text width =1em, font=\fontsize{8}{9}\selectfont] {抑菌效果};

        \draw[->,shift ={(0 ,-11.5)}] (5,-1) -- (5,-1.5);
        \draw[rounded corners,fill=green!10,shift ={(0 ,-13)}] (0,0) rectangle (10,-1) node [midway] {壳聚糖-辣椒果胶复合膜对冷藏草鱼片的防腐保鲜效果};

            \draw[dashed,fill=yellow!10,shift ={(1.8 ,-14.25)}] (0,0) rectangle (1,-3) node [midway,text width =1em, font=\fontsize{8}{9}\selectfont] {草鱼样品预处理};
            \draw[dashed,fill=yellow!10,shift ={(3.1 ,-14.25)}] (0,0) rectangle (1,-3) node [midway,text width =1em, font=\fontsize{8}{9}\selectfont] {鱼肉的感官评价};
            \draw[dashed,fill=yellow!10,shift ={(4.4 ,-14.25)}] (0,0) rectangle (1,-3) node [midway,text width =1em, font=\fontsize{8}{9}\selectfont] {理化指标变化分析};
            \draw[dashed,fill=yellow!10,shift ={(5.7 ,-14.25)}] (0,0) rectangle (1,-3) node [midway,text width =1em, font=\fontsize{8}{9}\selectfont] {微生物分析};
            \draw[dashed,fill=yellow!10,shift ={(7 ,-14.25)}] (0,0) rectangle (1,-3) node [midway,text width =1em, font=\fontsize{8}{9}\selectfont] {蛋白质变化分析};
    \end{tikzpicture}
\end{document}
